\chapter{The topological target}
\label{ch:topology}

Which topological object to compute is the central design decision, and the
requirements of \cref{ch:problem} settle it. This chapter records the reasoning,
including two conclusions that were wrong in earlier drafts and are corrected
here.

\section{Three candidates}

\begin{table}[htbp]
\centering
\caption{Topological summaries of a scalar field, and what periodicity costs
  each of them.}
\label{tab:candidates}
\begin{tabular}{p{0.20\textwidth}p{0.34\textwidth}p{0.34\textwidth}}
\toprule
object & shape & cost on a periodic domain \\
\midrule
join (merge) tree &
  always a tree, on a torus as much as anywhere &
  nothing \\
contour tree &
  equals the Reeb graph only if the domain is simply connected &
  not available: invalid \\
Reeb graph &
  carries loops &
  loop surgery~\cite{tierny2009loopsurgery} \\
\bottomrule
\end{tabular}
\end{table}

\section{What a Reeb loop is, and is not}

A cycle in the Reeb graph is a one-parameter family of contours that leaves a
level, travels, and returns. Cycles arise from non-trivial $1$-cycles in the
\emph{domain}, not from the shape of any one contour. On a $\dsp$-torus
$\betti{1}(\Tor{\dsp}) = \dsp$, so the Reeb graph of a triply periodic field can
carry up to three independent loops, each meaning a structure that closes on
itself through the wrap. On a non-periodic box $\betti{1} = 0$, the Reeb graph is
a tree, and it is the contour tree.

\paragraph{It does not give the genus of a level set.} A saddle that adds a
handle to a contour without merging or splitting it is a \emph{degree-two} node
in the Reeb graph, and degree-two nodes are conventionally dropped: they sit in
the middle of an arc and are invisible in the output. Handle counts, tunnels and
genus evolution are therefore not readable from a Reeb graph. An earlier draft
claimed otherwise; the correction matters, because it removes the main reason
one might have paid for loop surgery.

The right tool for genus is far cheaper. Extract the surface, split it into
connected components, and compute
\begin{equation}
  \eulerchi \;=\; V - E + F, \qquad g \;=\; \frac{2 - \eulerchi}{2}
  \label{eq:euler}
\end{equation}
per closed orientable component: exact, linear in the element count, no sweep.
Its one prerequisite is that the surface be watertight, which is what
\cref{ch:complex} delivers.

\paragraph{And $\eulerchi$ is the wrong invariant in two dimensions,} where every
closed curve has $\eulerchi = 0$ before and after any event of interest. The 2D
statement is the \emph{number} of boundary components together with their
homology class on the torus: a curve that wraps has no inside, so ``contained
within'' is conditional on being contractible. Winding numbers---net minimum-image
displacement around the loop, divided by the domain length---distinguish the two,
and non-contractible curves necessarily occur in pairs with opposite class,
which is itself a consistency check.

\section{The decision}
\label{sec:decision}

Cells are connected components of a super-level set, and their appearance,
growth and merging is a merge-tree story. Two consequences decide the matter.

\paragraph{Persistence is the robustness measure \ref{req:robust} asked for.}
Rather than choosing a threshold, compute all thresholds at once. Each cell then
carries a birth value $b$, a merge value $m$, and a persistence
\begin{equation}
  p \;=\; |b - m|,
  \label{eq:persistence}
\end{equation}
which is how far the knob must move to make that cell vanish. The reportable
quantity is $N(p)$, the number of cells with persistence above $p$. A plateau in
$N(p)$ is a defensible scale; a featureless continuum says no threshold is
defensible, which is itself the finding and shifts the question rather than
answering it badly.

\paragraph{Periodicity costs nothing.} A join tree is always a tree: components
of $\sublev{v}$ merge monotonically as $v$ decreases and can never form a cycle.
Loop surgery is required for the contour tree and the Reeb graph, neither of
which is needed. The doubly and triply periodic cases of \ref{req:periodic},
with arbitrary face placement, come free---periodicity enters only as adjacency
in the union--find, which \cref{ch:complex} already supplies. The most expensive
of the three candidates was the one not required.

\section{What the merge tree does not cover}

A crescent whose tips join, leaving one contour inside another, is a
$\betti{1}$ event: the region remains one connected component and the merge tree
is blind to it. Concretely, the region is a topological disk with one boundary
curve before the event and an annulus with two after, while
$\betti{0} = 1$ throughout and $\betti{1}$ goes $0 \to 1$. Capturing it requires
the per-component boundary invariants above---boundary-curve counts with homology
class in 2D, \cref{eq:euler} per surface component in 3D---which sit on the
segmentation and are small.

\section{One engine for two comparisons}
\label{sec:overlap}

Matching segmentations between consecutive snapshots by spatial overlap is the
same operation as matching a fuel-consumption segmentation against a temperature
segmentation on the \emph{same} snapshot. Built once it serves both: whether a
cell persists in time (\ref{req:track}), and whether two surrogates identify the
same cells (\ref{req:compare}). The second measures the open question of
\cref{ch:problem} directly---where surrogates agree the structure is
definition-independent, and where they diverge that is differential diffusion
expressed as topology rather than as a contour plot.

Segmentation written as a plotfile is what makes this practical: an integer label
per cell of the hierarchy feeds straight into \texttt{conditionalMean}, so
statistics of any field inside cells defined by any other field is a one-liner.
That is the common context the surrogates need in order to be compared at all,
and it is the presentation used in~\cite{bremer2010burning}.

\paragraph{A hypothesis, not an assertion.} A merge-tree branch that never
becomes an independent maximum---a saddle-separated shoulder of a larger
structure---is structurally distinct from one that is its own high-persistence
branch. If back-supported structures are the former and self-sustaining ones the
latter, then \cref{eq:persistence} measures the degree of independence. This
requires physics judgement and is recorded as a question.

\section{Flame surface thickness}

Broadening under turbulent strain follows from \cref{sec:carried} with no new
machinery: carry $|\nabla f|$ onto the surface as a sampled component, and the
local thickness is
\begin{equation}
  \delta \;\approx\; \frac{\Delta v}{|\nabla f|},
  \label{eq:thickness}
\end{equation}
aggregated over the surface as a distribution rather than a single number.
